% Extended Abstract — fNIRS 2025 / SfNIRS Submission Format
% processFNIRS2: A Modular Open-Source MATLAB Toolbox for fNIRS Analysis
% (Architecture-focused variant: ~75% architecture, ~25% FRESH validation)
%
% Max 1000 words (excl. title, authors, affiliations, figures, tables, references)
% Max 4 pages total.

\documentclass[11pt]{article}

% Page geometry — US Letter, top/bottom 0.8", left/right ~1"
\usepackage[letterpaper, top=0.8in, bottom=0.8in, left=1in, right=1in]{geometry}

% Fonts — Arial (via helvet)
\usepackage[utf8]{inputenc}
\usepackage[T1]{fontenc}
\usepackage{helvet}
\renewcommand{\familydefault}{\sfdefault}

% Figures and tables
\usepackage{graphicx}
\usepackage{booktabs}
\usepackage{caption}
\captionsetup{font=small, labelfont=it, textfont=it}

% Links
\usepackage{url}
\usepackage[colorlinks=true,linkcolor={blue},citecolor={blue},urlcolor={blue}]{hyperref}

% Spacing — single
\usepackage{setspace}
\setstretch{1.0}

% Color for headings
\usepackage{xcolor}
\definecolor{headblue}{HTML}{2E74B5}

% Section formatting — 12pt blue, not bold, left-aligned
\usepackage{titlesec}
\titleformat{\section}{\normalfont\fontsize{12}{14}\selectfont\color{headblue}}{}{0em}{}
\titlespacing*{\section}{0pt}{6pt}{6pt}

% Paragraph spacing
\setlength{\parindent}{0pt}
\setlength{\parskip}{6pt}

% Compact lists
\usepackage{enumitem}
\setlist{nosep}

\begin{document}

% ============================================================
% TITLE — 14pt Bold, Blue, Centered
% ============================================================
\begin{center}
{\fontsize{14}{17}\selectfont\bfseries\color{headblue}%
processFNIRS2: A Modular Open-Source MATLAB Toolbox for fNIRS Analysis, Validated on the FRESH Benchmark}

\vspace{6pt}

% AUTHORS — 11pt, Black, Centered
Adrian Curtin\textsuperscript{1}, Hasan Ayaz\textsuperscript{1,2,3,4}

\vspace{3pt}

% AFFILIATIONS — 11pt, Black, Centered
{\fontsize{11}{13}\selectfont
\textsuperscript{1}School of Biomedical Engineering, Science and Health Systems, Drexel University, Philadelphia, PA, USA\\
\textsuperscript{2}Department of Psychological and Brain Sciences, Drexel University, Philadelphia, PA, USA\\
\textsuperscript{3}Drexel Solutions Institute, Drexel University, Philadelphia, PA, USA\\
\textsuperscript{4}Center for Injury Research and Prevention, Children's Hospital of Philadelphia, Philadelphia, PA, USA}
\end{center}

% ============================================================
\section*{Synopsis}

We present processFNIRS2 (pf2), an open-source MATLAB toolbox built on a
modular, configuration-driven architecture spanning the full fNIRS workflow.
A two-layer design separates single-subject processing from group analysis: a
user-facing API package (\texttt{+pf2}) wraps an internal infrastructure layer
(\texttt{+pf2\_base}) and an interchangeable algorithm layer, while a dedicated
group-analysis package (\texttt{+exploreFNIRS}) handles LME modeling,
connectivity, and hyperscanning. Processing follows a three-stage pipeline
(raw~$\rightarrow$~optical density~$\rightarrow$~hemoglobin~$\rightarrow$
~filtered) whose steps are composable value-class pipelines built from
precomputed, immutable processing functions with dynamic argument injection. A
device-abstraction layer supports 16 hardware configurations through
declarative \texttt{.cfg} files, and a \texttt{ProcessingContext} object
enables globals-free, reproducible, parallel execution. To confirm numerical
correctness, we validated pf2 against three established toolboxes (Cedalion,
MNE-NIRS, AnalyzIR) on the FRESH benchmark, obtaining activation-map
correlations of \textit{r}~=~0.87--0.96.

% ============================================================
\section*{Motivation}

fNIRS analysis is fragmented across many toolboxes, each encoding different
processing choices, data structures, and interfaces. The FRESH study
(Yuecel et al., 2025) showed that 38 independent teams processing identical
fNIRS datasets produced hypothesis pass rates ranging from 25\% to 100\%,
underscoring how both implementation and analyst choices drive variability. A
toolbox that makes processing decisions explicit, composable, and reproducible
--- while remaining interoperable with community standards --- is therefore
valuable for transparent, multi-site work. processFNIRS2 (pf2) is designed
around these goals: a modular architecture in which every processing step is a
first-class, inspectable object; a configuration-driven core that supports
multiple devices without code changes; and an execution model that can bypass
global state entirely for reproducible and parallel analyses. Here we describe
the architecture and confirm, on the FRESH benchmark, that pf2 produces results
numerically consistent with three independently developed platforms.

% ============================================================
\section*{Architecture}

\textbf{Two-layer analysis model.} pf2 is organized into two layers with a
clean interface. Layer~1 (single-subject processing) imports raw device files
and converts them to hemoglobin biomarkers; Layer~2 (\texttt{+exploreFNIRS})
consumes collections of processed structs for group statistics, connectivity,
and hyperscanning. The contract between layers is a plain processed-fNIRS
struct, so subjects processed by any means can enter group analysis.

\textbf{Package layering.} Within Layer~1, a user-facing API package
(\texttt{+pf2}) provides import, export, data manipulation, probe/ROI, and
plotting namespaces. It wraps an internal infrastructure layer
(\texttt{+pf2\_base}) --- device loading, core fNIRS math, plotting style, and
wavelet transforms --- which in turn calls an interchangeable algorithm layer
(\texttt{/functions}) of stand-alone processing kernels (motion correction,
filtering, referencing). A configuration layer of declarative \texttt{.cfg}
files sits beneath all of it. A progressive-disclosure API mirrors this
hierarchy: parent-level calls (\texttt{pf2.import}) auto-select or prompt, while
explicit child calls (\texttt{pf2.import.importSNIRF}) give direct control.

\textbf{Three-stage pipeline.} Processing is factored into three stages ---
raw intensity~$\rightarrow$~optical density (motion correction, filtering,
referencing), optical density~$\rightarrow$~hemoglobin (modified Beer--Lambert
with selectable DPF modes, including age-dependent DPF after Scholkmann and
Wolf, 2013), and hemoglobin~$\rightarrow$~filtered (baseline correction,
filtering, artifact rejection, ROI averaging). Intermediate stages are
retained, making each transformation inspectable and reproducible.

\textbf{Composable pipeline value classes.} Each processing step is a
\texttt{PipelineFunction} --- an immutable value object that precomputes its
function handle and argument mapping. Ordered chains are \texttt{Pipeline}
objects (with \texttt{RawPipeline}/\texttt{OxyPipeline} specializations for
stages~1 and~3), supporting \texttt{add}, \texttt{insert}, \texttt{remove}, and
\texttt{setParam}, and round-tripping to/from named method configs. At runtime a
dynamic argument-injection mechanism supplies each kernel only the inputs it
declares (data matrix, sampling rate, channel mask, markers, auxiliary signals,
etc.), so generic functions compose without bespoke glue.

\textbf{Device abstraction.} A single codebase supports 16 fNIRS hardware
configurations through INI-style \texttt{.cfg} files describing wavelengths,
channel maps, and 2D/3D optode geometry. The \texttt{pf2.Device} value class
exposes clean accessors (wavelengths, MNI positions, source--detector
distances, channel tables) with a persistent cache, and is auto-attached to
imported data.

\textbf{Reproducible, globals-free execution.} For interactive use, settings
live in global state. For batch, parallel, and reproducible work, a
\texttt{ProcessingContext} object encapsulates all settings, methods, and the
device as locals; passing it to \texttt{processFNIRS2} bypasses globals
entirely --- they are neither read nor written --- and the context serializes
for exact reproduction. This enables isolated unit testing and one-context-per-
worker parallelism.

\textbf{Standards, visualization, and testing.} pf2 imports and exports SNIRF
and full BIDS-NIRS datasets, and exports a self-describing HDF5 tensor contract
for foundation-model workflows. Theme-aware plotting renders topographic maps,
3D cortical surfaces, connectomes, and DOT reconstructions headlessly (saved
figures always export on white). Correctness is guarded by a
\texttt{matlab.unittest} suite --- unit, integration, and quick-validation
tests --- backed by synthetic data generators for hemodynamics, noise, motion
artifacts, and physiology.

% ============================================================
\section*{Analysis Capabilities}

On this architecture, pf2 provides seven motion correction methods (TDDR
(Fishburn et al., 2019), wavelet (Molavi and Dumont, 2012), spline, SplineSG
(Jahani et al., 2018), sSMART (Ayaz et al., 2010), spline+wavelet hybrid, and
Takizawa rejection (Takizawa et al., 2008)); scalp-coupling-index quality
control; short-channel regression; block averaging and GLM with OLS or AR-IRLS
noise models; functional connectivity; hyperscanning with Hyper-Brain ICA
(Luo et al., 2025); graph-theoretic network metrics; and diffuse optical
tomography. Group analysis adds LME modeling with FDR correction and
publication-ready visualization.

% ============================================================
\section*{Validation}

To confirm numerical correctness, we implemented matched pipelines across pf2,
Cedalion (Python) (Middell et al., 2026), MNE-NIRS (Python) (Luke et al., 2021),
and AnalyzIR (MATLAB) (Santosa et al., 2018) on both FRESH datasets
(Yuecel et al., 2025): Dataset~I (auditory speech-noise, \textit{N}~=~17, 33
channels) and Dataset~II (motor finger-tapping, \textit{N}~=~10, 64
long-separation + 4 short-separation channels). All toolboxes used equivalent
steps (SNIRF import, SCI threshold 0.75, motion correction, Beer--Lambert
DPF~=~6, bandpass 0.01--0.5~Hz, short-channel regression, block averaging, and
AR-IRLS GLM). Twenty-two pf2 variants spanned five motion correction methods,
two noise models, and with/without short-channel regression and SCI QC.
Agreement used Pearson correlation of channel-wise activation, ICC(3,1)
(Shrout and Fleiss, 1979), S{\o}rensen--Dice similarity of significance maps,
and Bland--Altman analysis.

Activation maps were highly correlated across all toolbox pairs
(Figure~\ref{fig:correlation}): \textit{r}~=~0.87--0.96 for auditory block
averaging and 0.82--0.94 for standardized GLM, with Dice coefficients of
0.94--0.95 for motor significance maps. Raw ICC(3,1) was near zero (0.01--0.02)
because DPF/PPF conventions produce up to 10\textsuperscript{6}-fold magnitude
differences ($\mu$M vs. mol/L); after per-toolbox z-scoring, ICC rose to 0.86
(auditory) and 0.42 (motor). Individual-level motor pass rates were concordant
across all four toolboxes (Table~\ref{tab:motor}): combined contralateral
activation reached 80--100\% (block averaging) and 70--100\% (GLM), all at or
above the FRESH median and within the 27th--100th percentile of the FRESH team
distribution.

% ============================================================
% TABLE 1
% ============================================================
\begin{table}[ht]
\centering
\caption{Dataset~II individual-level motor hypothesis pass rates (\textit{X}/10 subjects). H1: left hand 2s FT activates contralateral PMC; H2: left + right 2s FT; H3: all sessions (2s+3s); H4: 3s > 2s duration effect. BA = block averaging, GLM = AR-IRLS.}
\label{tab:motor}
\small
\begin{tabular}{llcccc}
\toprule
Toolbox & Method & H1 & H2 & H3 & H4 \\
\midrule
pf2         & BA  & 4/10 & 9/10  & 9/10  & 1/10 \\
            & GLM & 6/10 & 8/10  & 9/10  & 2/10 \\
Cedalion    & BA  & 5/10 & 10/10 & 10/10 & 1/10 \\
            & GLM & 2/10 & 7/10  & 9/10  & 2/10 \\
MNE-NIRS    & BA  & 6/10 & 8/10  & 9/10  & 4/10 \\
            & GLM & 5/10 & 7/10  & 9/10  & 7/10 \\
AnalyzIR    & BA  & 4/10 & 8/10  & 10/10 & 6/10 \\
            & GLM & 8/10 & 10/10 & 10/10 & 0/10 \\
\midrule
FRESH median &    & 5/10 & 8/10  & 9/10  & 2/10 \\
\bottomrule
\end{tabular}
\end{table}

% ============================================================
% FIGURE 1
% ============================================================
\begin{figure}[ht]
\centering
\includegraphics[width=0.48\textwidth]{comparison/output/datasetI_correlation.png}
\caption{Inter-toolbox Pearson correlation of channel-wise group-level HbO activation (Dataset~I, block averaging). All pairwise correlations exceed \textit{r} = 0.75, with pf2 correlations ranging from 0.87 (MNE-NIRS) to 0.96 (AnalyzIR).}
\label{fig:correlation}
\end{figure}

% ============================================================
% FIGURE 2
% ============================================================
\begin{figure}[ht]
\centering
\includegraphics[width=0.48\textwidth]{comparison/output/datasetII_brain3d_topfront.png}
\caption{Group-level motor activation (\textit{t}-statistic) from GLM AR-IRLS with short-channel regression (Dataset~II), rendered with pf2's built-in headless 3D cortical visualization. Bilateral motor cortex activation is consistent with the finger-tapping paradigm. Sources (yellow) and detectors (red) shown.}
\label{fig:brain3d}
\end{figure}

% ============================================================
\section*{Discussion}

The architecture makes processing decisions explicit and composable:
pipelines are inspectable objects, devices are declarative configs, and a
\texttt{ProcessingContext} guarantees that a serialized analysis reproduces
byte-for-byte without global-state coupling. The FRESH validation confirms this
design is not only transparent but numerically faithful --- pf2 implements core
fNIRS algorithms (Beer--Lambert conversion, TDDR, bandpass filtering, GLM
estimation) consistently with three independently developed platforms. The
highest correlation was between pf2 and AnalyzIR (\textit{r}~=~0.96), both
MATLAB toolboxes sharing numerical libraries, yet pf2 also correlated strongly
with the Python toolboxes (\textit{r}~=~0.87--0.92), indicating the agreement is
not an artifact of shared implementation. The near-zero raw ICC highlights that
DPF/PPF and unit conventions --- not algorithmic disagreement --- dominate
apparent cross-toolbox differences; we recommend documenting these conventions
explicitly and using scale-invariant metrics for cross-platform comparison. The
concordant motor results reinforce the FRESH conclusion that analytical choices
drive more variability than toolbox implementation. We will extend this
comparison to connectivity and hyperscanning metrics at the conference.

% ============================================================
\section*{Availability}

processFNIRS2 is open-source (GPLv3) and released with documentation, runnable
tutorials, and the comparison harness used for this validation.

% ============================================================
\section*{Disclosures}

H. Ayaz was involved in the technology development of fNIR Devices, LLC and was thus offered a minor share in the startup firm. fNIR Devices, LLC manufactures optical brain imaging instruments and licensed IP and know-how from Drexel University.

% ============================================================
\section*{References}

{\small
Yuecel, M. A., et al. (2025). fNIRS Reproducibility and Estimation of Statistical power from a Harmonized study (FRESH). \textit{Commun. Biol.} doi: 10.1038/s42003-025-08412-1

Luke, R., Shader, M. J., McAlpine, D., Wiggins, I. M., Teng, J., Badcock, N. A., and Burnham, D. (2021). Analysis methods for measuring passive auditory fNIRS responses generated by a novel audiometry paradigm. \textit{Neurophotonics} 8, 025008.

Santosa, H., Zhai, X., Fishburn, F., and Huppert, T. (2018). The NIRS Brain AnalyzIR Toolbox. \textit{Algorithms} 11, 73.

Fishburn, F. A., Ludlum, R. S., Vaidya, C. J., and Medvedev, A. V. (2019). Temporal Derivative Distribution Repair (TDDR): a motion correction method for fNIRS. \textit{NeuroImage} 184, 171--179.

Shrout, P. E. and Fleiss, J. L. (1979). Intraclass correlations: uses in assessing rater reliability. \textit{Psychol. Bull.} 86, 420--428.

Middell, E., et al. (2026). Cedalion Tutorial: A Python-based framework for comprehensive analysis of multimodal fNIRS \& DOT from the lab to the everyday world. Preprint.

Molavi, B. and Dumont, G. A. (2012). Wavelet-based motion artifact removal for functional near-infrared spectroscopy. \textit{Physiol. Meas.} 33, 259--270.

Jahani, S., Setarehdan, S. K., Boas, D. A., and Y\"ucel, M. A. (2018). Motion artifact detection and correction in functional near-infrared spectroscopy: a new hybrid method based on spline interpolation method and Savitzky-Golay filtering. \textit{Neurophotonics} 5, 015003.

Ayaz, H., Izzetoglu, M., Shewokis, P. A., and Onaral, B. (2010). Sliding-window motion artifact rejection for Functional Near-Infrared Spectroscopy. \textit{Conf. Proc. IEEE Eng. Med. Biol. Soc.} 2010, 6567--6570.

Takizawa, R., Kasai, K., Kawakubo, Y., Marumo, K., Kawasaki, S., Yamasue, H., and Kato, N. (2008). Reduced frontopolar activation during verbal fluency task in schizophrenia: a multi-channel near-infrared spectroscopy study. \textit{Schizophr. Res.} 99, 250--262.

Scholkmann, F. and Wolf, M. (2013). General equation for the differential pathlength factor of the frontal human head depending on wavelength and age. \textit{J. Biomed. Opt.} 18, 105004. doi: 10.1117/1.JBO.18.10.105004

Luo, H., Cai, Y., Lin, X., and Duan, L. (2025). Hyper-brain independent component analysis (HB-ICA): an approach for detecting inter-brain networks from fNIRS-hyperscanning data. \textit{Biomed. Opt. Express} 16, 245--256.
}

\end{document}
